\documentclass[a4paper,12pt]{article}
\usepackage[utf8]{inputenc} % Pacote para acentuação
\usepackage[brazil]{babel} % Para quem vai escrever em português brasileiro
\usepackage[lmargin=3cm,tmargin=3cm,rmargin=2cm,bmargin=2cm]{geometry} %Formato que lembra a ABNT
\usepackage{amsmath,amsthm,amsfonts,amssymb,dsfont,mathtools} %pacotes matemáticos
\usepackage{graphicx}


\begin{document}

\begin{table}
\centering 
\begin{tabular}{l|c|c|c}\hline
  \textbf{Matérias}      & \textbf{Prova 1} & \textbf{Prova 2} & \textbf{Nota Final} \\\hline
  Cálculo 1              &  7,5 & 7,0 & $NFC_1$\\
  Álgebra Linear         &  8,5 & 6,0 & $NFA$\\
  Química 1              &  6,5 & 6,0 & $NFQ_1$\\
  Intro. às Engenharias  &  9,5 & 10,0& $NFI%\\
  
\end{tabular}\\
    \caption{relações de matéria}
\label{Tab1}
\end{table}

\begin{document}

\begin{table}
\centering 
\begin{tabular}{l|c|c|c}\hline
  \textbf{Matérias}      & \textbf{Prova 1} & \textbf{Prova 2} & \textbf{Nota Final} \\\hline
  Cálculo 1              &  7,5 & 7,0 & $NFC_1$\\
  Álgebra Linear         &  8,5 & 6,0 & $NFA$\\
  Química 1              &  6,5 & 6,0 & $NFQ_1$\\
  Intro. às Engenharias  &  9,5 & 10,0& $NFI%\\
  
\end{tabular}\\
    \caption{relações de reprovações}
    \label{tab2}
\end{table}
    Etou falando da tabela (ver tabela \ref{Tab1} e não da de reprovações(tabela \ref{tab2}
\end{document}
\begin{document}\begin{table}\centering \begin{tabular}{l|c|c|c}\hline  \textbf{Matérias}      & \textbf{Prova 1} & \textbf{Prova 2} & \textbf{Nota Final} \\\hline  Cálculo 1              &  7,5 & 7,0 & NFC_1\\  Álgebra Linear         &  8,5 & 6,0 & NFA\\  Química 1              &  6,5 & 6,0 & NFQ_1\\  Intro. às Engenharias  &  9,5 & 10,0& NFI\\\end{tabular}\\\label{Tabela de Notas e Avaliações}\end{table}\end{document}
